% Generated draft: requires mathematical and source review.
\documentclass[11pt]{article}
\usepackage[margin=25mm]{geometry}
\usepackage{fontspec}
\setmainfont{Latin Modern Roman}
\usepackage{amsmath,amssymb,mathtools,unicode-math}
\setmathfont{Latin Modern Math}
\usepackage{enumitem,needspace,xurl,hyperref}
\setlength{\parindent}{0pt}
\setlength{\parskip}{5pt}
\setlength{\emergencystretch}{3em}
\newcommand{\N}{\mathbb{N}}
\newcommand{\Z}{\mathbb{Z}}
\newcommand{\Q}{\mathbb{Q}}
\newcommand{\R}{\mathbb{R}}
\newcommand{\C}{\mathbb{C}}
\newcommand{\F}{\mathbb{F}}
\newcommand{\A}{\mathbb{A}}
\newcommand{\PP}{\mathbb{P}}
\newcommand{\E}{\mathbb{E}}
\newcommand{\eps}{\varepsilon}
\newcommand{\dd}{\,\mathrm{d}}
\begin{document}
\section*{UnsolvedMath ID: 30004567}
\textbf{Source problem code:} catalogue entry
\subsection{Definitions and mathematical statement}
Let \(K\) denote a knot or link type, and write \(\operatorname{Cr}(K)\) for its crossing number, meaning the minimum number of crossings among regular diagrams representing \(K\). The supplied record concerns folded ribbon realizations \(K_w\), with width parameter \(w>0\), and their folded ribbonlength \(\operatorname{Rib}(K_w)\). It does not define admissible folded ribbons, their length normalization, allowed overlaps or degeneracies, ribbon equivalence, or framing conventions. Thus a complete geometric definition cannot be recovered from this record alone.

For a fixed choice of those conventions, let \(\mathcal A(K)\) denote the admissible folded ribbon realizations of \(K\), and, whenever this class is nonempty, put
\[
R_{\mathcal A}(K):=\inf_{K_w\in\mathcal A(K)}\operatorname{Rib}(K_w).
\]
This notation makes the dependence on the unspecified admissibility conventions explicit; it does not supply those conventions.

\begin{enumerate}
\item Determine lower bounds for folded ribbonlength of knots and links. An invariant lower bound would take the form
\[
\operatorname{Rib}(K_w)\geq B(K)
\qquad\text{for every }K_w\in\mathcal A(K),
\]
equivalently \(R_{\mathcal A}(K)\geq B(K)\). The source does not specify the desired function \(B\), its dependence on crossing number or other invariants, or an intended sharpness criterion.

\item Investigate whether folded ribbon realizations can be constructed with a crossing-number upper bound of exponent \(p\) satisfying \(1\leq p<3/2\). If the intended bound is universal over the knot and link types under consideration, its explicit quantifiers are
\[
\begin{gathered}
\exists p\in[1,3/2)\ \exists C>0\ \exists N_0\geq1\ \forall K:\\
\operatorname{Cr}(K)\geq N_0
\ \Longrightarrow\
\exists K_w\in\mathcal A(K)\quad
\operatorname{Rib}(K_w)\leq C\,\operatorname{Cr}(K)^p.
\end{gathered}
\]
Here \(p,C,N_0\) are independent of \(K\). The source does not settle whether this uniform scope is intended for all knots and links, for knots alone, or for particular families. This asymptotic formulation leaves types with crossing number below \(N_0\) outside the stated bound.

\item Investigate families with superlinear folded ribbonlength in crossing number. If this refers to the infimum over all admissible realizations, a precise formulation is to seek types \((K_n)_{n\geq1}\) with nonempty admissible classes such that
\[
\operatorname{Cr}(K_n)\longrightarrow\infty,
\qquad
\frac{R_{\mathcal A}(K_n)}{\operatorname{Cr}(K_n)}\longrightarrow\infty.
\]
The source does not specify whether this limiting formulation, a particular superlinear rate, or merely growth of selected realizations is intended. These are different targets. The upper-bound and superlinear-family questions are retained as separate questions; the source gives no precise logical alternative between them.

\item Determine how folded ribbonlength depends on ribbon equivalence and framing, including which admissible classes and infima are to be compared when either is fixed or varied. The equivalence relation, framing data, and requested comparison are not supplied, so this part remains an explicitly unresolved formulation task.
\end{enumerate}
\subsection{Short English statement}
The problem asks for lower bounds and improved crossing-number upper bounds for folded ribbonlength, for families whose optimal folded ribbonlength grows superlinearly, and for the effects of ribbon equivalence and framing. The displayed formulations are conditional on geometric conventions absent from the supplied record. This is a statement-only draft requiring human mathematical review.
\subsection{Sources}
\begin{itemize}
\item Supplied local UnsolvedMath registry record, permanent ID \(30004567\), title \emph{Folded Ribbonlength Bounds and Knot Complexity}, source problem code \texttt{catalogue entry}, in the Oberwolfach Reports Open Problems collection. No individual proposer is identified.
\item Supplied report citation: \emph{Mini-Workshop: Nonlocal Analysis and the Geometry of Embeddings (2021)}, Oberwolfach Reports, pp.\ 1823--1856, DOI \texttt{10.4171/owr/2020/37}. \url{https://doi.org/10.4171/owr/2020/37}. The record describes Oberwolfach Reports as documenting mathematical workshops and their open-problem discussions. Its proposed year and year tag are 2020; the supplied citation labels the workshop 2021. These details are retained without reconciliation.
\item Supplied literature reference: Denne, folded ribbonlength. \url{https://arxiv.org/abs/2010.03611}.
\item Supplied literature reference: Kim et al., ribbonlength of two-bridge knots. \url{https://arxiv.org/abs/2208.03669}.
\item Supplied literature reference: Denne and Patterson, ribbonlength of torus and twist knots. \url{https://arxiv.org/abs/2509.18370}.
\item Literature assessment checked at \texttt{2026-08-22}: the supplied status is \emph{partially solved}. The supplied assessment reports general folded-ribbon upper bounds, linear bounds for important knot families, and recent bounded examples, while describing the proposed universal polynomial or superlinear characterization as unsettled. This is supplied dated curation derived from the OWR record, not an independently verified literature assessment; the status and cited claims require independent verification.
\item The registry labels the statement \texttt{corrected\_verified}, but its verification note says that only formula delimiters, notation, grammar, and typography were repaired, with mathematical content and scope preserved. This label does not establish mathematical verification of the formulations above.
\end{itemize}
\end{document}
